\documentclass[11pt]{article}

\usepackage[margin=1in]{geometry}
\usepackage[T1]{fontenc}
\usepackage[utf8]{inputenc}
\usepackage{lmodern}
\usepackage{amsmath,amssymb,bm}
\usepackage{booktabs}
\usepackage{enumitem}
\usepackage{hyperref}

\hypersetup{
  colorlinks=true,
  linkcolor=blue,
  urlcolor=blue,
  pdftitle={Theoretical Foundations for VLEO Slew Maneuver}
}

\setlength{\parindent}{0pt}
\setlength{\parskip}{6pt}

\newcommand{\vect}[1]{\boldsymbol{#1}}
\newcommand{\transpose}{\mathsf{T}}
\newcommand{\norm}[1]{\left\lVert #1 \right\rVert}
\newcommand{\diag}{\operatorname{diag}}

\title{Theoretical Foundations\\\large VLEO Slew Maneuver}
\author{}
\date{}

\begin{document}

\maketitle
\tableofcontents
\newpage

\section{Coordinate Transformations}

\subsection{Quaternion to Direction Cosine Matrix}

A quaternion $\vect{q} = [q_0\; q_1\; q_2\; q_3]^\transpose$ using the
scalar-first convention can be converted to a direction cosine matrix
through

\begin{equation}
\mathbf{C}(\vect{q}) =
\begin{bmatrix}
q_0^2 + q_1^2 - q_2^2 - q_3^2 & 2(q_1 q_2 + q_0 q_3) & 2(q_1 q_3 - q_0 q_2) \\
2(q_1 q_2 - q_0 q_3) & q_0^2 - q_1^2 + q_2^2 - q_3^2 & 2(q_2 q_3 + q_0 q_1) \\
2(q_1 q_3 + q_0 q_2) & 2(q_2 q_3 - q_0 q_1) & q_0^2 - q_1^2 - q_2^2 + q_3^2
\end{bmatrix}.
\end{equation}

\textbf{Source.} Wertz, J. R., \emph{Spacecraft Attitude Determination and
Control}, 1978, Eq. 12-12.

\textbf{Implemented in.} Aerospace Toolbox \texttt{quat2dcm}, used in
\texttt{src/Sat\_template.m}, \texttt{src/sat\_template\_gui.m},
\texttt{src/gui/AttitudeAnimator.m},
\texttt{src/gui/AttitudeAnimatorGUI.m}, and
\texttt{src/computeAeroForces.m}.

\subsection{Direction Cosine Matrix to Quaternion}

The inverse transformation uses Shepperd's method for numerical
stability:

\begin{enumerate}[label=\arabic*.]
\item Compute the matrix trace and diagonal elements.
\item Find the maximum of $\{\operatorname{tr}(\mathbf{C}), C_{11}, C_{22}, C_{33}\}$.
\item Use the corresponding extraction formula to avoid division by a
small number.
\end{enumerate}

\textbf{Source.} Shepperd, S. W., ``Quaternion from Rotation Matrix,''
\emph{Journal of Guidance and Control}, Vol. 1, No. 3, 1978.

\textbf{Implemented in.} Aerospace Toolbox \texttt{dcm2quat}, used in
\texttt{src/Reference\_Trajectory\_Calculation.m} and
\texttt{src/gui/runSimulation.m}.

\subsection{Quaternion and Principal Axis-Angle}

Quaternion from axis-angle:

\begin{equation}
q_0 = \cos\left(\frac{\theta}{2}\right), \qquad
q_1 = e_1 \sin\left(\frac{\theta}{2}\right), \qquad
q_2 = e_2 \sin\left(\frac{\theta}{2}\right), \qquad
q_3 = e_3 \sin\left(\frac{\theta}{2}\right).
\end{equation}

Here $\vect{e} = [e_1\; e_2\; e_3]^\transpose$ is the unit rotation axis
and $\theta$ is the rotation angle.

Axis-angle from quaternion:

\begin{equation}
\theta = 2 \cos^{-1}(q_0),
\end{equation}

with

\begin{equation}
\vect{e} =
\begin{cases}
\dfrac{1}{\sin(\theta/2)} [q_1\; q_2\; q_3]^\transpose, & \theta \neq 0, \\
[1\; 0\; 0]^\transpose, & \theta = 0.
\end{cases}
\end{equation}

\textbf{Source.} Euler's rotation theorem; Kuipers, J.,
\emph{Quaternions and Rotation Sequences}, 1999.

\textbf{Implemented in.} No current \texttt{src/} implementation. If
principal-axis-angle support is needed again, add it back as a documented
custom repository function.

\subsection{Fundamental Rotation Matrices}

Elementary rotations about the principal axes are

\begin{align}
\mathbf{R}_1(\theta) &=
\begin{bmatrix}
1 & 0 & 0 \\
0 & \cos\theta & \sin\theta \\
0 & -\sin\theta & \cos\theta
\end{bmatrix}, \\
\mathbf{R}_2(\theta) &=
\begin{bmatrix}
\cos\theta & 0 & -\sin\theta \\
0 & 1 & 0 \\
\sin\theta & 0 & \cos\theta
\end{bmatrix}, \\
\mathbf{R}_3(\theta) &=
\begin{bmatrix}
\cos\theta & \sin\theta & 0 \\
-\sin\theta & \cos\theta & 0 \\
0 & 0 & 1
\end{bmatrix}.
\end{align}

\textbf{Source.} Schaub, H. and Junkins, J., \emph{Analytical Mechanics
of Space Systems}, 4th ed., 2018, Ch. 3.

\textbf{Implemented in.} No dedicated \texttt{src/} helper. Current
\texttt{src/} code uses Aerospace Toolbox rotation and conversion
utilities instead of a local elementary-rotation wrapper.

\section{Orbital Mechanics}

\subsection{Orbital Elements from Position and Velocity}

Given inertial position $\vect{r}$ and velocity $\vect{v}$, the classical
orbital elements are computed through

\begin{enumerate}[label=\arabic*.]
\item Specific angular momentum:
\begin{equation}
\vect{h} = \vect{r} \times \vect{v}.
\end{equation}

\item Node vector:
\begin{equation}
\vect{n} = [0\; 0\; 1]^\transpose \times \vect{h}.
\end{equation}

\item Eccentricity vector:
\begin{equation}
\vect{e} = \frac{\left(\norm{\vect{v}}^2 - \mu / \norm{\vect{r}}\right)\vect{r} - (\vect{r} \cdot \vect{v})\vect{v}}{\mu}.
\end{equation}

\item Specific orbital energy:
\begin{equation}
\epsilon = \frac{\norm{\vect{v}}^2}{2} - \frac{\mu}{\norm{\vect{r}}}.
\end{equation}

\item Semi-major axis:
\begin{equation}
a = -\frac{\mu}{2\epsilon}.
\end{equation}

\item Inclination:
\begin{equation}
i = \cos^{-1}\left(\frac{h_z}{\norm{\vect{h}}}\right).
\end{equation}

\item Right ascension of the ascending node:
\begin{equation}
\Omega = \cos^{-1}\left(\frac{n_x}{\norm{\vect{n}}}\right),
\end{equation}
with the quadrant selected from the sign of $n_y$.

\item Argument of periapsis:
\begin{equation}
\omega = \cos^{-1}\left(\frac{\vect{n} \cdot \vect{e}}{\norm{\vect{n}}\norm{\vect{e}}}\right),
\end{equation}
with the quadrant selected from the sign of $e_z$.

\item True anomaly:
\begin{equation}
f = \cos^{-1}\left(\frac{\vect{e} \cdot \vect{r}}{\norm{\vect{e}}\norm{\vect{r}}}\right),
\end{equation}
with the quadrant selected from the sign of $\vect{r} \cdot \vect{v}$.
\end{enumerate}

\textbf{Special cases.}
\begin{itemize}
\item Circular orbits ($e \approx 0$): use the argument of latitude
$u = \omega + f$.
\item Equatorial orbits ($i \approx 0$): use the longitude of periapsis
$\varpi = \Omega + \omega$.
\end{itemize}

\textbf{Source.} Vallado, D. A., \emph{Fundamentals of Astrodynamics and
Applications}, 4th ed., 2013, Algorithm 9.

\textbf{Implemented in.} Aerospace Toolbox \texttt{ijk2keplerian} when
state-to-elements conversion is needed.

\subsection{Position and Velocity from Orbital Elements}

First compute the position and velocity in the perifocal frame:

\begin{align}
\vect{r}_{\mathrm{PQW}} &= \frac{a(1-e^2)}{1 + e\cos f}
\begin{bmatrix}
\cos f \\
\sin f \\
0
\end{bmatrix}, \\
\vect{v}_{\mathrm{PQW}} &= \sqrt{\frac{\mu}{a(1-e^2)}}
\begin{bmatrix}
-\sin f \\
e + \cos f \\
0
\end{bmatrix}.
\end{align}

Then transform to the ECI frame with

\begin{align}
\mathbf{R} &= \mathbf{R}_3(-\Omega)\mathbf{R}_1(-i)\mathbf{R}_3(-\omega), \\
\vect{r}_{\mathrm{ECI}} &= \mathbf{R}\vect{r}_{\mathrm{PQW}}, \\
\vect{v}_{\mathrm{ECI}} &= \mathbf{R}\vect{v}_{\mathrm{PQW}}.
\end{align}

\textbf{Source.} Vallado, D. A., \emph{Fundamentals of Astrodynamics and
Applications}, 4th ed., 2013, Algorithm 10.

\textbf{Implemented in.} Aerospace Toolbox \texttt{keplerian2ijk}, used
in \texttt{src/Reference\_Trajectory\_Calculation.m},
\texttt{src/gui/runSimulation.m}, and
\texttt{src/gui/AttitudeAnimator.m}.

\section{Attitude Dynamics}

\subsection{Quaternion Kinematics}

The time derivative of the quaternion representing body orientation is

\begin{equation}
\dot{\vect{q}} = \frac{1}{2}\,\mathbf{\Omega}(\vect{\omega})\vect{q},
\end{equation}

where the quaternion rate matrix is

\begin{equation}
\mathbf{\Omega}(\vect{\omega}) =
\begin{bmatrix}
0 & -\omega_1 & -\omega_2 & -\omega_3 \\
\omega_1 & 0 & \omega_3 & -\omega_2 \\
\omega_2 & -\omega_3 & 0 & \omega_1 \\
\omega_3 & \omega_2 & -\omega_1 & 0
\end{bmatrix}.
\end{equation}

Here $\vect{\omega} = [\omega_1\; \omega_2\; \omega_3]^\transpose$ is
the angular velocity of the body frame relative to the inertial frame,
expressed in body coordinates.

\textbf{Source.} Wertz, J. R., \emph{Spacecraft Attitude Determination and
Control}, 1978, Section 12.1.

\textbf{Implemented in.} \texttt{src/Sat\_template.m}.

\subsection{Euler's Rotational Equations}

For a rigid body with inertia tensor $\mathbf{J}$,

\begin{equation}
\mathbf{J}\dot{\vect{\omega}} = -\vect{\omega} \times (\mathbf{J}\vect{\omega}) + \vect{T}_{\mathrm{ext}}.
\end{equation}

For a diagonal inertia tensor $\mathbf{J} = \diag([J_x, J_y, J_z])$,

\begin{align}
J_x \dot{\omega}_x &= (J_y - J_z)\omega_y\omega_z + T_x, \\
J_y \dot{\omega}_y &= (J_z - J_x)\omega_z\omega_x + T_y, \\
J_z \dot{\omega}_z &= (J_x - J_y)\omega_x\omega_y + T_z.
\end{align}

\textbf{Source.} Schaub, H. and Junkins, J., \emph{Analytical Mechanics
of Space Systems}, 4th ed., 2018, Ch. 4.

\textbf{Implemented in.} \texttt{src/Sat\_template.m}.

\subsection{PD Attitude Control Law}

Quaternion-based feedback control is written as

\begin{equation}
\vect{T}_{\mathrm{control}} = -\mathbf{K}_p\vect{q}_{\mathrm{error},v} - \mathbf{K}_d\vect{\omega}_{\mathrm{error}},
\end{equation}

where

\begin{itemize}
\item $\vect{q}_{\mathrm{error},v}$ is the vector part of the attitude
error quaternion,
\item $\vect{\omega}_{\mathrm{error}}$ is the angular velocity error,
\item $\mathbf{K}_p$ is the proportional gain matrix, and
\item $\mathbf{K}_d$ is the derivative gain matrix.
\end{itemize}

\textbf{Source.} Wie, B., \emph{Space Vehicle Dynamics and Control}, 2nd
ed., 2008, Ch. 7.

\textbf{Implemented in.} \texttt{src/controllers/control\_torques.m}.

\section{Perturbations}

\subsection{J2 Perturbation (Earth Oblateness)}

The acceleration due to Earth's $J_2$ zonal harmonic is

\begin{equation}
\vect{a}_{J_2} = \frac{3}{2} J_2 \frac{\mu}{r^2} \left(\frac{R_E}{r}\right)^2
\begin{bmatrix}
\left(\dfrac{x}{r}\right)\left(5\left(\dfrac{z}{r}\right)^2 - 1\right) \\
\left(\dfrac{y}{r}\right)\left(5\left(\dfrac{z}{r}\right)^2 - 1\right) \\
\left(\dfrac{z}{r}\right)\left(5\left(\dfrac{z}{r}\right)^2 - 3\right)
\end{bmatrix}.
\end{equation}

Here $J_2 = 1.08263 \times 10^{-3}$ is dimensionless,
$R_E = 6378137\,\mathrm{m}$ is the Earth's equatorial radius,
$\mu = 3.986004418 \times 10^{14}\,\mathrm{m}^3/\mathrm{s}^2$ is the
Earth gravitational parameter, and $r = \norm{\vect{r}}$ is the distance
from the Earth's center.

\textbf{Source.} Vallado, D. A., \emph{Fundamentals of Astrodynamics and
Applications}, 4th ed., 2013, Eq. 8-25.

\textbf{Implemented in.} Aerospace Toolbox \texttt{gravityzonal}, used in
\texttt{src/Sat\_template.m} and \texttt{src/sat\_template\_gui.m}.

\subsection{J3 Perturbation}

The acceleration due to Earth's $J_3$ zonal harmonic is

\begin{equation}
\vect{a}_{J_3} = \frac{1}{2} J_3 \frac{\mu}{r^2} \left(\frac{R_E}{r}\right)^3
\begin{bmatrix}
\left(\dfrac{x}{r}\right)\left(10\left(\dfrac{z}{r}\right)^3 - 15\left(\dfrac{z}{r}\right)\right) \\
\left(\dfrac{y}{r}\right)\left(10\left(\dfrac{z}{r}\right)^3 - 15\left(\dfrac{z}{r}\right)\right) \\
4\left(\dfrac{z}{r}\right)^3 - 3\left(\dfrac{z}{r}\right)
\end{bmatrix},
\end{equation}

with $J_3 = -2.5327 \times 10^{-6}$.

\textbf{Source.} Vallado, D. A., \emph{Fundamentals of Astrodynamics and
Applications}, 4th ed., 2013, Eq. 8-26.

\textbf{Implemented in.} Aerospace Toolbox \texttt{gravityzonal}, used in
\texttt{src/Sat\_template.m} and \texttt{src/sat\_template\_gui.m}.

\subsection{Atmospheric Drag in VLEO}

The drag acceleration model is

\begin{equation}
\vect{a}_{\mathrm{drag}} = -\frac{1}{2} \rho C_d \frac{A}{m}
\norm{\vect{v}_{\mathrm{rel}}}\vect{v}_{\mathrm{rel}}.
\end{equation}

Here $\rho$ is the atmospheric density, $C_d$ is the drag coefficient,
$A$ is the cross-sectional area, $m$ is the spacecraft mass, and
$\vect{v}_{\mathrm{rel}}$ is the velocity relative to the atmosphere.

\textbf{Atmospheric models.}
\begin{itemize}
\item Exponential atmosphere (simple)
\item NRLMSISE-00 (high fidelity)
\item JB2008 (operational)
\end{itemize}

\textbf{Source.} Vallado, D. A., \emph{Fundamentals of Astrodynamics and
Applications}, 4th ed., 2013, Ch. 8.6.

\textbf{Implemented in.} \texttt{src/computeAeroForces.m}, which evaluates
NRLMSISE-00 using geodetic coordinates resolved from either direct
latitude/longitude inputs or the spacecraft ECI position and simulation
time. The resulting forces and moments are applied in
\texttt{src/Sat\_template.m}, \texttt{src/sat\_template\_gui.m}, and
\texttt{src/gui/runSimulation.m}.

\subsection{Solar Radiation Pressure}

The solar radiation pressure acceleration model is

\begin{equation}
\vect{a}_{\mathrm{SRP}} = -P_{\mathrm{sr}} C_r \frac{A}{m}
\frac{\vect{r}_{\mathrm{sun}}}{\norm{\vect{r}_{\mathrm{sun}}}}.
\end{equation}

Here $P_{\mathrm{sr}} \approx 4.56 \times 10^{-6}\,\mathrm{N/m}^2$ at
1 AU, $C_r$ is the reflectivity coefficient, $A$ is the illuminated area,
and $\vect{r}_{\mathrm{sun}}$ is the position vector from the spacecraft to
the Sun.

\textbf{Source.} Montenbruck, O. and Gill, E., \emph{Satellite Orbits},
2000, Ch. 3.4.

\textbf{Implemented in.} \texttt{src/computeAeroForces.m}, which computes
panel-wise solar coefficients and returns total solar force and moment
contributions for the active spacecraft mesh.

\section{Propulsion}

\subsection{Ion Thruster Fundamentals}

The thrust relation is

\begin{equation}
F = \dot{m} v_e = \dot{m} g_0 I_{sp},
\end{equation}

where $F$ is thrust, $\dot{m}$ is mass flow rate, $v_e$ is exhaust
velocity, $g_0 = 9.80665\,\mathrm{m/s}^2$ is standard gravity, and
$I_{sp}$ is specific impulse.

The corresponding power-thrust relation is

\begin{equation}
P = \frac{1}{2} \dot{m} v_e^2 \frac{1}{\eta},
\end{equation}

where $\eta$ is the thruster efficiency.

\textbf{Source.} Goebel, D. and Katz, I., \emph{Fundamentals of Electric
Propulsion: Ion and Hall Thrusters}, JPL, 2008.

\textbf{Implemented in.} TBD.

\subsection{Slew Maneuver Analysis}

For a rest-to-rest slew maneuver of angle $\theta$, the minimum-time
bang-bang control result is

\begin{equation}
t_{\min} = 2\sqrt{\frac{\theta J}{T_{\max}}},
\end{equation}

and a smooth minimum-energy form is written as

\begin{equation}
E_{\min} = \frac{4}{3} \frac{J\theta^2}{t^3}.
\end{equation}

Here $J$ is the moment of inertia about the rotation axis and
$T_{\max}$ is the maximum available torque.

\textbf{Source.} Wie, B., \emph{Space Vehicle Dynamics and Control}, 2nd
ed., 2008, Ch. 7.

\textbf{Implemented in.} TBD.

\section{References}

\subsection{Textbooks}

\begin{enumerate}[label=\arabic*.]
\item Vallado, D. A. (2013). \emph{Fundamentals of Astrodynamics and
Applications}, 4th ed. Microcosm Press.
\item Wertz, J. R. (1978). \emph{Spacecraft Attitude Determination and
Control}. Kluwer Academic Publishers.
\item Schaub, H. and Junkins, J. (2018). \emph{Analytical Mechanics of
Space Systems}, 4th ed. AIAA.
\item Wie, B. (2008). \emph{Space Vehicle Dynamics and Control}, 2nd ed.
AIAA.
\item Montenbruck, O. and Gill, E. (2000). \emph{Satellite Orbits:
Models, Methods, and Applications}. Springer.
\item Kuipers, J. (1999). \emph{Quaternions and Rotation Sequences}.
Princeton University Press.
\item Goebel, D. and Katz, I. (2008). \emph{Fundamentals of Electric
Propulsion: Ion and Hall Thrusters}. JPL/Wiley.
\end{enumerate}

\subsection{Papers}

\begin{enumerate}[label=\arabic*.]
\setcounter{enumi}{7}
\item Shepperd, S. W. (1978). ``Quaternion from Rotation Matrix.''
\emph{Journal of Guidance and Control}, Vol. 1, No. 3.
\end{enumerate}

\subsection{Standards}

\begin{enumerate}[label=\arabic*.]
\setcounter{enumi}{8}
\item IERS Conventions (2010). IERS Technical Note No. 36.
\item WGS84 (2014). NGA.STND.0036\_1.0.0\_WGS84.
\end{enumerate}

\section*{Document History}
\addcontentsline{toc}{section}{Document History}

\begin{tabular}{@{}lll@{}}
\toprule
Date & Author & Changes \\
\midrule
2025-12-30 & Team & Initial template with existing implementations \\
\bottomrule
\end{tabular}

\end{document}
